\section{结论与后续工作}

\subsection{结论}

本阶段我完成了电磁多极展开从理论到数值的完整闭环：

\begin{enumerate}
  \item \textbf{理论上}，手推了从 Maxwell 方程到 Grahn 2012 映射关系、再到 Alaee 2018 精确多极展开（Table 1/Table 2）的完整矢量球谐链，并理解了 FC2015/FC2017 关于精确偶极矩与 toroidal 非独立性的论证；推导经 scipy 独立数值验证与逐条退化检查。
  \item \textbf{数值上}，复现了四篇论文的核心结果：球体的 Table 2--Mie 四通道 $<0.03\%$（金球）、Grahn 双路径映射误差低至 $10^{-7}$、双金盘的 COMSOL 频域真解谱（ED 主导 + MD/EQ 磁共振）。
  \item \textbf{验证上}，三条物理判据（多近似互验、极限退化、能量守恒/光学定理）在球体与 Grahn 映射上全部通过，在双盘上能量守恒与多近似（FEM/BEM）通过、极限退化待补。
\end{enumerate}

我的核心信心来自\textbf{三条相互独立的物理验证}，而不是「曲线看起来像」：任何一条不过，都说明算错了，而非「近似误差」。

\subsection{后续工作}

\begin{itemize}
  \item \textbf{BEM 加表面阻抗（IBC）}：当前边界元是理想导体（PEC）近似，丢失了金的损耗与 MD/EQ 磁共振；加表面阻抗 $Z_s=Z_0/(n+ik)$ 可恢复，并有望把 MD/EQ 的网格收敛推到 $1\%$ 以内（表面型方法内存省，可跑更细网格）。
  \item \textbf{Fig.~3 的极限退化}：双盘无解析长波目标，需构造耦合偶极子半解析模型作为长波退化的参考，或明确说明「球体已验证、双盘无解析长波解」。
  \item \textbf{高阶多极 $l>4$ 逐阶解析}：当前只给出聚合上界，逐阶解析可进一步确认多极分解的完整性。
  \item \textbf{Fig.~2 的严格不确定性量化}：读图数据退役后需重新预注册电子数据口径的 UQ 契约。
\end{itemize}
